\section{Round-ten reconstruction: interior source saddles and a regenerative mixed-shell coefficient}
\label{sec:r10-b1}

Let
\[
 C(z)=\left(1,v,\frac{|v|^2}{2},\chi_1(z),\ldots,\chi_d(z)\right)
\]
be the complete initial constraint mark.  We impose the model's necessary
nondegeneracy in its intrinsic form: the affine span of \(C\) on the positive
activity support is the full declared constraint affine space.  Equivalently,
no nonzero linear combination of the added \(\chi_j\) is almost surely a
combination of particle number, momentum and energy.  The target \(a\) lies in
a compact subset of the relative interior of the mean image.  These are
exactly the conditions under which a finite source saddle and a
subexponential local coefficient can exist.

\subsection{The full finite-volume saddle}

For a bounded path/contact source \(H\), define
\[
 Q_\varepsilon(H,\lambda)=\mu_\varepsilon^{-1}
 \log\mathbb E^{\rm gc}_\varepsilon
 e^{\mu_\varepsilon\langle H,X_\varepsilon\rangle+
       \lambda\cdot C_{\rm tot}}.
\]

\begin{lemma}[Typical-sector covariance]
\label{lem:r10-b1-covariance}
On compact source and multiplier charts,
\[
 cI\le D^2_{\lambda\lambda}Q_\varepsilon(H,\lambda)\le CI
\]
on the declared affine constraint space, uniformly in small \(\varepsilon\).
The lower bound comes from linearly many particles in disjoint activity boxes,
not from an event with a fixed particle number.
\end{lemma}

\begin{proof}
Choose finitely many positive-activity phase boxes whose centered mark vectors
span the affine constraint space.  Divide the physical torus into
\(c\mu_\varepsilon\) mutually separated insertion cells.  The hard-core
cluster expansion gives a conditional occupancy probability bounded away
from zero and one for a positive fraction of those cells, uniformly on the
source chart.  Conditional variance in the spanning mark directions is
therefore at least \(c\mu_\varepsilon\); division by the pressure speed gives
the lower bound.  The upper bound follows from the uniform second connected
cumulant estimate.  Exponentially rare fixed-particle sectors are never used.
\end{proof}

\begin{corollary}[Exact source-dependent saddle]
\label{cor:r10-b1-saddle}
For every regular \((H,a)\) there is a unique finite-volume multiplier
\(\lambda_{\varepsilon,H,a}\) satisfying
\[
 D_\lambda Q_\varepsilon(H,\lambda_{\varepsilon,H,a})=a.
\]
It converges locally uniformly, with source derivatives, to the unique
limiting multiplier \(\lambda_{H,a}\).
\end{corollary}

\begin{proof}
Apply the quantitative inverse-function theorem using
Lemma~\ref{lem:r10-b1-covariance}, normal convergence of the full connected
pressure, and compactness of the interior target set.
\end{proof}

\subsection{Regenerative smoothing of the complete compound law}

A single mark may live on a curved lower-dimensional surface.  We therefore
use blocks.  Fix \(m_0\) disjoint insertion cells and call a block good when
exactly one particle lies in each of a selected set of cells and their phase
points fall in product patches on which the sum map
\[
 (z_1,\ldots,z_{m_0})\longmapsto\sum_{j=1}^{m_0}C(z_j)
\]
has Jacobian rank equal to the constraint dimension.  The affine-span
condition and the constant-rank theorem provide finitely many such patches.

\begin{lemma}[Good-block minorization with exponentially small complement]
\label{lem:r10-b1-goodblocks}
The configuration contains \(c\mu_\varepsilon\) disjoint candidate blocks.
On the full source chart, each block is good with conditional probability at
least \(p_*>0\), up to a summable cluster correction.  Hence
\[
 \mathbb P(\hbox{fewer than }c_*\mu_\varepsilon
 \hbox{ good blocks})\le e^{-c\mu_\varepsilon}.
\]
Conditioned on the exterior variables, the sum contributed by every good block
has a \(W^{s,1}\) density with uniform norm for any fixed \(s\) after enlarging
\(m_0\) finitely.
\end{lemma}

\begin{proof}
Separated insertion cells interact only through a convergent hard-core
polymer correction.  The finite-energy insertion estimate gives the uniform
conditional lower probability.  Exponential Chebyshev bounds for the block
indicators, with the cluster correction absorbed into the exponent, give the
first statement.  On a good product patch, coarea gives a density with a
uniform first Sobolev norm.  Convolution of finitely many independent
conditional block contributions increases the Sobolev order; bounded
conditional density ratios from the polymer expansion preserve the uniform
bounds.  Thus every typical configuration, not merely one local component,
contains linearly many smoothing blocks.
\end{proof}

\begin{theorem}[Full mixed lattice--continuous characteristic estimate]
\label{thm:r10-b1-characteristic}
Under the exact tilted law at
\(\lambda_{\varepsilon,H,a}\), the centered joint characteristic function of
particle number and continuous constraints satisfies, uniformly in the source
chart,
\[
 |\varphi_\varepsilon(t,u)|\le
 \begin{cases}
 e^{-c\mu_\varepsilon(t^2+|u|^2)},&|t|+|u|\le\delta,\\
 e^{-c\mu_\varepsilon},&\delta\le |t|\le\pi,
                            \ |u|\le R,\\
 e^{-c\mu_\varepsilon}+C_s(1+|u|)^{-s},&|u|>R,
 \end{cases}
\]
for arbitrarily large fixed \(s\).
\end{theorem}

\begin{proof}
The central estimate is the Taylor expansion at the exact finite saddle and
Lemma~\ref{lem:r10-b1-covariance}.  Away from the number lattice origin, a
positive fraction of insertion cells yields the usual span-one factor
\(|1+ze^{it}|/(1+z)<1\).  For large continuous frequency, split according to
the event in Lemma~\ref{lem:r10-b1-goodblocks}.  Its complement is
exponentially small.  On the good event, integrate by parts across the
linearly many Sobolev block densities; using only a fixed number of blocks
already gives the prescribed polynomial order, while the remaining blocks
preserve the bound.  Thus no positive-mass singular remainder survives in the
high-frequency estimate.
\end{proof}

\subsection{The admissible weighted shell}

Let particle number be fixed at its lattice target.  Let the continuous shell
be
\[
 |Y-a'|\le\delta_\varepsilon,
 \qquad
 \delta_\varepsilon\downarrow0,
 \qquad
 \sqrt{\mu_\varepsilon}\,\delta_\varepsilon\longrightarrow\infty.
\]
The target is the mean at the exact finite saddle.

\begin{theorem}[Uniform exact-saddle shell coefficient]
\label{thm:r10-b1-shell}
Uniformly on compact regular source/target charts,
\[
 \mathbb P_{H,\lambda_{\varepsilon,H,a},\varepsilon}
 \{N=N_\varepsilon,\ |Y-a'|\le\delta_\varepsilon\}
 =\mu_\varepsilon^{-1/2}(c_{H,a}+o(1)),
\]
where \(0<c\le c_{H,a}\le C\).  In particular the logarithm of the shell
probability is \(o(\mu_\varepsilon)\).
\end{theorem}

\begin{proof}
Use Fourier inversion in the exact lattice coordinate and smooth upper/lower
approximations of the continuous box.  The central region gives the joint
Gaussian coefficient at the exact saddle.  Since the continuous window
contains a diverging number of central-limit standard deviations, its
conditional Gaussian mass tends to one, while the exact number coefficient is
of order \(\mu_\varepsilon^{-1/2}\).  The two minor arcs and the large
continuous-frequency region are negligible by
Theorem~\ref{thm:r10-b1-characteristic}.  Uniform covariance bounds give the
positive coefficient bounds.
\end{proof}

\begin{theorem}[Source-dependent microcanonical pressure and prepared LDP]
\label{thm:r10-b1-main}
Conditioning on the admissible shell yields
\[
 Q^{\rm mc,a}(H)=
 Q(H,\lambda_{H,a})-\lambda_{H,a}\cdot a
 -Q(0,\lambda_{0,a})+\lambda_{0,a}\cdot a.
\]
The initial empirical measure has the constrained grand-canonical entropy
rate, and its Gaussian covariance is the Schur complement of the full
constraint covariance.  The theorem applies jointly to the B2 particle and
actual-contact source class.
\end{theorem}

\begin{proof}
Change measure in the numerator to the exact finite saddle.  On the shell the
multiplier factor equals
\(e^{-\mu_\varepsilon\lambda_{\varepsilon,H,a}\cdot a+o(\mu_\varepsilon)}\),
and Theorem~\ref{thm:r10-b1-shell} makes the residual shell cost
subexponential.  Repeat with \(H=0\) for the denominator and pass to the
normal pressure limit.  Convex duality gives the constrained initial rate.
Differentiating the saddle equation gives the covariance Schur complement.
\end{proof}
